\documentclass[10pt]{article}
\usepackage[margin=0.9in]{geometry}
\usepackage{amsmath,amssymb,amsthm,enumitem,hyperref}
\hypersetup{hidelinks,pdftitle={Degree and square-root energy on the Boolean cube},
 pdfauthor={ChatGPT-6 Astra; Kaarel Hanni; Dmitry Vaintrob}}
\newtheorem{theorem}{Theorem}
\newtheorem{lemma}[theorem]{Lemma}
\newtheorem{corollary}[theorem]{Corollary}
\newcommand{\E}{\mathbb E}
\newcommand{\Prb}{\mathbb P}
\newcommand{\En}{\mathcal E}
\newcommand{\F}{\mathbb F}
\newcommand{\sgn}{\operatorname{sgn}}
\title{Degree and square-root energy on the Boolean cube:\\
equivalent positive estimates and a signed obstruction}
\author{ChatGPT-6 Astra \and Kaarel H\"anni \and Dmitry Vaintrob\thanks{Authors are listed in alphabetical order.}}
\date{October 16, 2026}
\begin{document}
\maketitle
\begin{abstract}
For nonnegative Boolean-cube polynomials, we show that any universal
square-root energy bound which vanishes with the degree-to-dimension
ratio implies a linear degree bound. The same assertion is equivalent
to approximation, in relative $L^1$ error, by sums of squares of
polynomials of degree proportional to the original degree. One fixed
relative approximation error strictly below one suffices. The proof
uses exact linear-code embeddings and tensor powers. For arbitrary
signed polynomials the analogous square-root statement is false, even
with a global unit absolute margin and with all Fourier coefficients
on one level. The signed obstruction persists at any fixed prediction accuracy below
one for full parity, without a bound on the predictor's source cost.
These are auxiliary analytic results. The positive inequality itself,
and the intended natural-scale posterior sample
complexity bound for sparse parity, remain open.
\end{abstract}

\section{Definitions and the positive equivalence}
All expectations on $\{-1,1\}^k$ use the uniform probability measure.
Write $\chi_S(x)=\prod_{i\in S}x_i$ and
\[
 \Delta\chi_S=|S|\chi_S,\qquad
 \En(q)=\E q\Delta q=\sum_S |S|\widehat q(S)^2.
\]
Degree always means exact multilinear Walsh degree. For $P\ge0$,
the quantity considered below is $\En(\sqrt P)$, and
$\|\sqrt P\|_2^2=\E P$. Zero polynomials and constants are harmless
and may be omitted from statements involving their degree in a denominator.

\begin{theorem}\label{thm:equivalence}
The following universal assertions are equivalent. Each is quantified
over every dimension $k$ and every nonzero nonnegative polynomial $P$
on the $k$-cube of positive degree $r$.
\begin{enumerate}[label=(\roman*),leftmargin=*]
\item There is a finite absolute constant $C$ such that
\[
 \En(\sqrt P)\le Cr\E P.
\]
\item There is $\omega:[0,1]\to[0,\infty)$ with $\omega(t)\to0$
as $t\downarrow0$ such that
\[
 \En(\sqrt P)\le k\omega(r/k)\E P.
\]
\item For each $\varepsilon>0$ there is a finite $C_\varepsilon$ such that
\[
 \sum_{|S|>C_\varepsilon r}\widehat{\sqrt P}(S)^2
       \le\varepsilon\E P.
\]
\item For each $\varepsilon>0$ there is a finite $D_\varepsilon$ such that
there are real polynomials $Q_j$ with
\[
 \deg Q_j\le D_\varepsilon r,\qquad
 \E\left|P-\sum_jQ_j^2\right|\le\varepsilon\E P.
\]
\end{enumerate}
Moreover, (iii) for one fixed $\varepsilon<1$ already implies (i).
Likewise, (iv) for one fixed relative error $\varepsilon<1$ already
implies (i). If (iv) holds for all errors, its approximant can always
be chosen to be a single square, with a possibly different degree constant.
\end{theorem}

The theorem is an equivalence of existence statements. It does not
establish any of them or assert a sharp constant. In particular,
a weak degree-fraction estimate on the full positive cone is not an
easier existence target than the linear estimate.

\section{An exact noise embedding}
We identify cube coordinates with $\F_2$ bits when applying a binary
linear map. This identification does not change the Walsh notation.

\begin{lemma}\label{lem:code}
Let $\theta\in(0,1/2)$ be rational. For every $k$ there are an integer
$K$ and a surjective binary linear map $A:\F_2^K\to\F_2^k$, with at least
one zero column, such that
the lift $\widetilde P(Y)=P(AY)$ preserves the distribution of $P$ and
sends a Walsh character of degree $j$ to a character of degree
\[
 d_\theta(j)=\frac K2\bigl(1-(1-2\theta)^j\bigr)\le K\theta j.
\]
For every real function $q$ on the original cube,
\begin{equation}\label{eq:energy-code}
 \frac{\En(q\circ A)}K
  =\frac12\sum_S\bigl(1-(1-2\theta)^{|S|}\bigr)\widehat q(S)^2.
\end{equation}
\end{lemma}
\begin{proof}
Let $\mu_\theta$ be the law on $\F_2^k$ with independent
Bernoulli$(\theta)$ coordinates. Choose a common denominator $K$ for
all its masses, and form a $k$-by-$K$ matrix with exactly
$K\mu_\theta(v)$ copies of each $v$ as columns. The zero column is
allowed and merely creates an irrelevant physical variable. Every
standard basis vector occurs, so the matrix has full row rank and
$AY$ is uniform when $Y$ is uniform. The lift of $\chi_S$ has degree
$K\Prb_{v\sim\mu_\theta}(v\cdot S=1)=d_\theta(|S|)$.
Distinct characters remain distinct by full row rank. Parseval's
identity now gives \eqref{eq:energy-code}.
\end{proof}

\section{Proof of the positive equivalence}
First, (i) implies (ii) and (iii), using Markov's inequality for the
nonnegative Fourier masses of $\sqrt P$.

Suppose (ii). We may replace $\omega$ near zero by a nondecreasing
envelope, which still tends to zero. Apply Lemma~\ref{lem:code} to $P$.
Since $\deg\widetilde P/K\le\theta r$, it gives
\begin{equation}\label{eq:weak-noise}
 \frac12\sum_S\bigl(1-(1-2\theta)^{|S|}\bigr)
               \widehat{\sqrt P}(S)^2
       \le\omega(\theta r)\E P.
\end{equation}
For rational $c\in(0,1/2)$ take $\theta=c/r$. If $|S|>r/c$,
then $(1-2\theta)^{|S|}<e^{-2}$. Consequently
\[
 \sum_{|S|>r/c}\widehat{\sqrt P}(S)^2
 \le\frac{2\omega(c)}{1-e^{-2}}\E P.
\]
Letting $c$ be a sufficiently small fixed rational proves (iii).

Next suppose (iii) holds for one fixed $\varepsilon<1$, with constant
$C_\varepsilon=C$. Normalize $\E P=1$ and put $q=\sqrt P$.
The masses $\widehat q(S)^2$ define a probability distribution;
let $J$ be the degree $|S|$ under this distribution. Then
$\E J=\En(q)$. For the tensor power
\[
 P_m(x^{(1)},\ldots,x^{(m)})=\prod_{a=1}^mP(x^{(a)}),
\]
degree is $mr$, and the Fourier degree distribution of $\sqrt{P_m}$
is the law of $J_1+\cdots+J_m$ for independent copies of $J$.
Thus $\Prb(J_1+\cdots+J_m>Cmr)\le\varepsilon$ for every $m$.
If $\E J>Cr$, the weak law of large numbers makes this probability
tend to one. Therefore $\En(\sqrt P)\le Cr$, proving (i).

For (iii)$\Rightarrow$(iv), let $Q$ be the Fourier projection of
$q=\sqrt P$ to degrees at most $C_\delta r$. Cauchy--Schwarz gives
\[
 \E|P-Q^2|\le\|q-Q\|_2(\|q\|_2+\|Q\|_2)
       \le2\sqrt\delta\,\E P.
\]
Choose $\delta=(\varepsilon/2)^2$.

Conversely, suppose that for one fixed $\delta<1$ there is
$R=\sum_jQ_j^2$ with $\deg Q_j\le Dr$ and
$\E|P-R|\le\delta\E P$. The Euclidean norm is Lipschitz, so
on every edge
\[
 (\sqrt{R(x)}-\sqrt{R(y)})^2
       \le\sum_j(Q_j(x)-Q_j(y))^2.
\]
It follows that $\En(\sqrt R)\le Dr\E R\le Dr(1+\delta)\E P$.
Also $\|\sqrt P-\sqrt R\|_2^2\le\E|P-R|$. Hence
\[
 \frac{\|\Pi_{>Cr}\sqrt P\|_2}{\sqrt{\E P}}
    \le\sqrt\delta+\sqrt{\frac{D(1+\delta)}C}.
\]
For a sufficiently large fixed $C$, the right side is strictly below
one. We have obtained a fixed version of (iii), which already implies
(i). This proves all assertions of Theorem~\ref{thm:equivalence}.

There is also a direct constant conversion. In \eqref{eq:weak-noise},
apply the embedding to $P_m$ with $\theta=c/(mr)$. The left side,
after dividing by $\E P_m$, tends to
\[
 \frac12\left[1-\exp\left(-\frac{2c\En(\sqrt P)}{r\E P}\right)\right].
\]
Thus whenever $\omega(c)<1/2$,
\begin{equation}\label{eq:constant}
 \frac{\En(\sqrt P)}{r\E P}
       \le\frac{-\log(1-2\omega(c))}{2c}.
\end{equation}
The embeddings are finite and exact, although their dimensions can be
very large.

In fact, one fixed improvement on the trivial bound suffices: if there
are $c>0$ and $\eta<1/2$ such that
\[
 \En(\sqrt P)\le\eta k\E P
 \quad\text{whenever }P\ge0\text{ and }\deg P\le ck,
\]
then the same tensor-code argument proves (i). For any positive rational
$c'<\min(c,1/2)$ it gives
$C=-\log(1-2\eta)/(2c')$. Here the trivial universal bound is
$\En(\sqrt P)\le k\E P/2$, obtained by expanding each edge square
and using $\sqrt{P(x)P(y)}\ge0$.

\section{A signed obstruction, including a global absolute margin}
For a real function $P$ define its signed square root
\[
 M(P)=\sgn(P)\sqrt{|P|},\qquad \|M(P)\|_2^2=\E|P|.
\]

\begin{theorem}\label{thm:signed}
There is a sequence of real polynomials $G_j$ on cubes of dimensions
$N_j$, with degrees $r_j$, such that
\[
 |G_j(x)|\ge1\quad\hbox{for every }x,\qquad
 \frac{r_j}{N_j}\longrightarrow0,\qquad
 \frac{\En(M(G_j))}{N_j\E|G_j|}\ge\frac{1-e^{-1}}4.
\]
Each $G_j$ has all its nonzero Walsh coefficients on a single level.
Consequently no universal bound
$\En(M(P))\le k\omega(\deg(P)/k)\E|P|$ with $\omega(t)\to0$
can hold for arbitrary signed polynomials, even under the displayed
global absolute-margin condition.
\end{theorem}

\begin{proof}
For odd $m$, set $S_m=\sum_{i=1}^m x_i$ and
$a_m=\En(M(S_m))/\E|S_m|$. We first verify
\begin{equation}\label{eq:affine-log}
 a_m=\frac18\log m+O(1).
\end{equation}
Let $T$ be a sum of $m-1$ independent signs, and put $p_0=\Prb(T=0)$.
Coordinate symmetry gives
\[
 \En(M(S_m))=\frac m4\E\bigl(M(T+1)-M(T-1)\bigr)^2,
       \qquad \E|S_m|=mp_0.
\]
The squared increment is $4$ at $T=0$. For positive even $t\ge2$,
it equals
\[
 d(t)=\frac4{(\sqrt{t+1}+\sqrt{t-1})^2}
       =\frac1t+O(t^{-3}),
\]
and the expression is even in $t$. The local binomial estimate
$\Prb(T=t)/p_0=1+O((t^2+1)/m)$, uniformly for even
$2\le t\le\sqrt m$, yields
\[
 2\sum_{\substack{2\le t\le\sqrt m\\t\text{ even}}}
       \frac{\Prb(T=t)}{p_0}d(t)
       =\frac12\log m+O(1).
\]
The local-estimate errors and the $O(t^{-3})$ terms are summable with
total $O(1)$. For $t>\sqrt m$, use $d(t)\le C/\sqrt m$ and
$p_0\asymp m^{-1/2}$ to bound the remaining normalized contribution
by $O(1)$. Dividing the energy identity by $mp_0$ proves
\eqref{eq:affine-log}.

Fix odd $m$ and form the $\ell$-fold tensor product
$F_\ell=\prod_{b=1}^\ell S_m(x^{(b)})$. It is homogeneous of
Walsh degree $\ell$, and $|F_\ell|\ge1$. The signed square root
tensorizes. If $Z$ has the Walsh degree distribution of
$M(S_m)/\sqrt{\E|S_m|}$, then $\E Z=a_m$, and the corresponding
degree distribution for $M(F_\ell)$ is $Z_1+\cdots+Z_\ell$.

Choose rational $\theta_\ell\in(0,1/2)$ with
$\ell\theta_\ell\to1/(2a_m)$ and apply Lemma~\ref{lem:code}
to obtain $G_\ell=F_\ell\circ A$ on a $K_\ell$-cube. The lift
preserves $|G_\ell|\ge1$, and \eqref{eq:energy-code} gives
\[
 \frac{\En(M(G_\ell))}{K_\ell\E|G_\ell|}
   =\frac12\E\left[1-(1-2\theta_\ell)^{Z_1+\cdots+Z_\ell}\right]
   \longrightarrow\frac{1-e^{-1}}2.
\]
Here the law of large numbers and bounded convergence justify the
limit. Homogeneity and the exact character-degree formula give
\[
 \frac{\deg G_\ell}{K_\ell}
 =\frac12\bigl(1-(1-2\theta_\ell)^\ell\bigr)
 \longrightarrow\frac12(1-e^{-1/a_m})\le\frac1{2a_m}.
\]
For each increasing odd $m$, choose $\ell$ sufficiently large to make
the energy ratio at least $(1-e^{-1})/4$ and the degree fraction at
most $1/a_m$. Equation~\eqref{eq:affine-log} then gives the required
diagonal sequence. All nonzero characters of each $G_\ell$ have the
same physical degree, so the final homogeneity assertion also follows.
\end{proof}

The global margin in Theorem~\ref{thm:signed} is absolute.
Indeed $\E|G_\ell|=(\E|S_m|)^\ell$ can be enormous, so the theorem
does not assert a unit margin after normalization by $\E|G_\ell|$
or by a neural-network source cost. Nor is agreement with any specified
target parity imposed. The result therefore obstructs a universal
analytic shortcut, not the desired learning lower bound.

It also does not contradict a bound for the different nonlinear
Rayleigh functional $\E\sgn(P)\Delta P$. For the homogeneous
polynomials just constructed, this functional is exactly
$\deg(P)\E|P|$, while the signed-square-root energy has the behavior
in Theorem~\ref{thm:signed}. Finally, $|S_m|$ itself does not have
degree one; the signed counterexample makes no claim against the
nonnegative-polynomial conjecture in Theorem~\ref{thm:equivalence}.

\section{Fixed parity accuracy does not remove the signed obstruction}
The preceding counterexamples need not predict parity. A selector
construction adds arbitrary fixed accuracy while retaining the energy
obstruction. The absence of a source-cost bound remains essential to
the scope of the conclusion.

\begin{corollary}\label{cor:accuracy}
Fix $\alpha<1$. There is a sequence of real polynomials $P_j$ on
$n_j$-cubes such that
\[
 |P_j|\ge1\text{ everywhere},\qquad
 \frac{\deg P_j}{n_j}\longrightarrow0,\qquad
 \Prb\left(\sgn(P_j)=\prod_{i=1}^{n_j}x_i\right)\ge\alpha,
\]
and
\[
 \frac{\En(M(P_j))}{n_j\E|P_j|}\ge\frac{1-e^{-1}}{16}.
\]
No bound on actual neural-network source cost or posterior cell energy
is asserted.
\end{corollary}
\begin{proof}
Let $B$ on a $K$-cube be a member of the sequence constructed in
Theorem~\ref{thm:signed}, and put $b_0=(1-e^{-1})/4$. Its code has
zero columns by Lemma~\ref{lem:code}, so $B$ ignores at least one
physical bit. Hence $\sgn B$ agrees with full parity $f_K$ on exactly
half of the cube.

For any fixed $\varepsilon>0$ there is a polynomial $G$ of degree
$O_\varepsilon(\sqrt K)$ with $|G|\ge1$ everywhere and
$\Prb(\sgn G=f_K)\ge1-\varepsilon$. To see this, let
$h(y)=(K-\sum_i y_i)/2$ be Hamming weight. Choose integers $a\le b$
with $\Prb(a\le h\le b)\ge1-\varepsilon$ and
$b-a\le\sqrt{K/\varepsilon}+2$, using Chebyshev's inequality and
rounding. The polynomial
\[
 G_0(h)=(-1)^b\prod_{j=a}^{b-1}(h-j-1/2)
\]
has sign $(-1)^h$ for integer $a\le h\le b$, and has no zeros at
any integer. Divide by its positive minimum absolute value on
$\{0,\ldots,K\}$ to obtain $G$.

Add $t$ selector bits $z$, where $t$ is fixed, and write
$H(z)=\prod_{i=1}^t(1+z_i)/2$, $\chi(z)=\prod_{i=1}^t z_i$,
and $p=2^{-t}$. For $A\ge1$ set
\[
 P(z,y)=H(z)A B(y)+(1-H(z))\chi(z)G(y).
\]
This has global absolute margin one, degree at most
$t+\max(\deg B,\deg G)=o(K)$, and full-parity accuracy at least
\[
 p/2+(1-p)(1-\varepsilon).
\]
Choose the fixed $t$ and $\varepsilon$ so this is at least $\alpha$.
The old-coordinate energies and nonnegativity of selector-coordinate
energies give
\[
 \En(M(P))\ge pA\En(M(B)),\qquad
 \E|P|=pA\E|B|+(1-p)\E|G|.
\]
Choose $A\ge(1-p)\E|G|/(p\E|B|)$. For $K\ge t$ this yields
\[
 \frac{\En(M(P))}{(K+t)\E|P|}
 \ge\frac{b_0K}{2(K+t)}\ge\frac{b_0}{4}.
\]
Taking the sequence of $B$'s proves the corollary.
\end{proof}

The multiplier $A$ separates prediction accuracy from normalized
signed-root energy. The corollary does not estimate the repair cost of
this selector family: neither large amplitude nor large signed-root
energy alone supplies that estimate. Excluding these examples from a
model with uniformly bounded repair cost would require a separate
geometric proof for that model. The analytic obstruction established
here therefore leaves calibration for accurate heads with bounded
actual source cost, and the natural-scale learning target, open.
\end{document}
